\chapter{Summary, Conclusion, and Recommendations}\label{ch:5}

\section{Summary of Findings}\label{sec:5-findings}

This study developed and evaluated \textit{StatiCalcs}, an interactive web-based learning tool with integrated calculators designed for the subject \textit{Statics of Rigid Bodies}. The study determined the level of perception of Bachelor of Science in Civil Engineering students and instructors of MSU--General Santos regarding the platform’s usability, accessibility, and satisfaction, as well as their level of preference for different learning resources.

Based on the findings, the following were obtained:

\begin{enumerate}
    \item \textbf{Level of Perception of Students Regarding StatiCalcs}

    The student respondents demonstrated a positive level of perception regarding StatiCalcs in all three evaluated areas. In terms of usability, the platform obtained an overall weighted mean of 4.66, interpreted as Strongly Agree, indicating that students perceived the website as easy to navigate, efficient, and user-friendly. In terms of accessibility, StatiCalcs obtained an overall weighted mean of 4.65, interpreted as Strongly Agree, showing that the website was perceived as readable, functional across devices, and easy to understand. In terms of satisfaction, the platform achieved an overall weighted mean of 4.63, interpreted as Strongly Agree, indicating a high level of satisfaction among students regarding the website’s usefulness for practice and problem solving.

    \item \textbf{Level of Perception of Instructors Regarding StatiCalcs}

    The instructor respondents likewise demonstrated a positive level of perception regarding StatiCalcs. In terms of usability, the platform obtained an overall weighted mean of 4.45, interpreted as Strongly Agree, indicating that instructors perceived the website as easy to navigate and usable. In terms of accessibility, StatiCalcs obtained an overall weighted mean of 4.29, interpreted as Strongly Agree, reflecting favorable evaluations regarding readability, device compatibility, and accessibility. In terms of satisfaction, the platform achieved an overall weighted mean of 4.18, interpreted as Agree, showing positive instructor evaluations regarding the usefulness and reliability of the platform.

    \item \textbf{Significant Difference Between the Level of Perception of Students and Instructors Regarding StatiCalcs}

    The independent samples \textit{t}-test results revealed that there was no significant difference between the level of perception of students and instructors regarding StatiCalcs in terms of usability, accessibility, and satisfaction. The computed \textit{p}-values for usability (0.19), accessibility (0.22), and satisfaction (0.06) were all greater than the 0.05 level of significance. Therefore, the null hypothesis was accepted, indicating that both respondent groups shared similar levels of perception toward the developed web-based learning tool.

    \item \textbf{Level of Preference for Learning Resources in Studying Statics of Rigid Bodies}

    Among student respondents, \textit{StatiCalcs} ranked first with a weighted mean of 4.55, interpreted as Most Preferred, indicating a strong preference for interactive and technology-supported learning resources. Textbooks, YouTube videos, and modules were likewise highly preferred.

    Among instructor respondents, textbooks ranked first with a weighted mean of 4.45, interpreted as Most Preferred, while \textit{StatiCalcs} ranked third with a weighted mean of 3.73, interpreted as Preferred. These findings suggest that instructors continue to value traditional instructional materials while also recognizing the usefulness of digital learning tools such as \textit{StatiCalcs}.

    \item \textbf{Respondents' Feedback Regarding StatiCalcs}

    Although not part of the main research questions, the respondents generally provided positive feedback regarding \textit{StatiCalcs}, describing the platform as helpful, user-friendly, convenient, and beneficial for learning \textit{Statics of Rigid Bodies}. Several recommendations for improvement were also identified, including the addition of more examples, tutorials, visualization features, improved mobile accessibility, and expanded topic coverage.    

\end{enumerate}

%====================================================================

\section{Conclusions}\label{sec:5-conclusions}

Based on the findings of the study, the following conclusions were drawn:

\begin{enumerate}

\item The developed web-based learning tool, \textit{StatiCalcs}, was positively evaluated by student respondents in terms of usability, accessibility, and satisfaction. The findings indicate that the platform effectively provided a user-friendly, accessible, and satisfactory supplementary learning environment for studying \textit{Statics of Rigid Bodies}.

\item The instructor respondents likewise demonstrated a positive level of perception toward \textit{StatiCalcs} in terms of usability, accessibility, and satisfaction. This suggests that the platform possesses educational value and potential usefulness as a supplementary teaching and learning tool for engineering instruction.

\item There was no significant difference between the level of perception of students and instructors regarding \textit{StatiCalcs} in terms of usability, accessibility, and satisfaction. Therefore, the platform was similarly accepted and positively evaluated by both groups of respondents.

\item Student respondents identified \textit{StatiCalcs} as their most preferred learning resource in studying \textit{Statics of Rigid Bodies}, while instructors continued to favor textbooks but still recognized \textit{StatiCalcs} as a preferred supplementary learning tool. This indicates that digital learning platforms can effectively complement traditional instructional resources in engineering education.

\item The positive feedback, suggestions, and recommendations of the respondents indicate that \textit{StatiCalcs} has strong potential for further enhancement, wider implementation, and possible expansion into additional engineering subjects and functionalities.

\end{enumerate}

\section{Recommendations}\label{sec:5-recommendations}

Based on the findings, conclusions, and respondents’ feedback, the following recommendations are proposed:

\begin{enumerate}

\item Future developers may improve the learning support features of \textit{StatiCalcs} by incorporating more examples, practice problems, diagrams, tutorials, and clearer step-by-step explanations. These additions may further strengthen conceptual understanding and improve the overall learning experience of users.

\item Future versions of \textit{StatiCalcs} may enhance system functionality and accessibility by improving mobile responsiveness, navigation features, formula visibility, input validation, visualization tools, and support for additional units of measurement. Future developers may also consider incorporating offline functionality or limited offline access to selected features to improve usability in environments with unstable or limited internet connectivity.

\item Future developers may consider transforming \textit{StatiCalcs} into a dedicated mobile application to improve accessibility and convenience for users. This may allow students and instructors to access the platform more easily across different devices and learning environments.

\item Future development of the platform may expand its scope by incorporating additional engineering topics and subjects, such as \textit{Dynamics of Rigid Bodies} and \textit{Mechanics of Deformable Bodies}, to broaden the educational applicability of the system for engineering students.

\item Instructors may consider integrating \textit{StatiCalcs} into classroom instruction, review sessions, demonstrations, and practice-solving activities to further support the teaching and learning process in \textit{Statics of Rigid Bodies}. This may enhance student engagement, reinforce engineering concepts, and supplement traditional instructional materials.

\item Future studies may provide respondents with a longer exposure period to the platform before conducting evaluation and data gathering. Allowing students and instructors to use the system for a longer duration, such as throughout an academic term or semester, may provide more comprehensive feedback regarding its usability, accessibility, satisfaction, and educational value.

\item Future researchers may evaluate the effect of \textit{StatiCalcs} on students’ academic performance, problem-solving ability, and conceptual understanding in \textit{Statics of Rigid Bodies}. This may determine whether the platform contributes not only to positive levels of perception but also to measurable learning outcomes such as quiz scores, examination performance, and analytical skills.

\item Future researchers may conduct similar studies involving students from other engineering programs that take \textit{Statics of Rigid Bodies}, and not only Bachelor of Science in Civil Engineering students. This may help determine the applicability, effectiveness, and acceptance of \textit{StatiCalcs} across a wider engineering student population.

\end{enumerate}